In dimension one, within the standing BFJ hypotheses, the chain
stabilizes on any strictly semistable model \(\mathcal X\)
determining both \(\theta\) and the real model obstacle \(u\).
Indeed, \cref{lem:ch23-psh-vertex-chord} replaces each arbitrary
\(\theta\)-psh competitor below \(u\) by a larger real model
competitor determined on that same model.  Consequently
\[
                    P^{\mathcal X}_\theta(u)=P_\theta(u).
\]
The largest element in the fixed-model competitor set exists here:
the set of its vertex vectors is nonempty, bounded above by the
obstacle, closed, and closed under coordinatewise maxima; the
finite-coordinate supremum argument in the proof of
\cref{thm:ch23-model-orthogonality-curve} therefore applies.
Nonemptiness follows by applying the chord lemma to a global
competitor below \(u\), whose existence in this BFJ ample-class
regime was proved in \cref{cor:ch23-continuous-finite}.
If a finite extension is needed to obtain the model, this statement
is applied to the base-changed data; a conclusion back over the
original field requires the corresponding descent statement.

The preceding equality concerns the full BFJ class and real model
functions in this chapter.  In a different field regime,
\cref{thm:ch29-real-class-curve-model-envelope} instead proves a
theorem about rational-model envelopes.  One cannot substitute
that theorem for the chord lemma or identify the two competitor
classes without proof.  The one-dimensional mechanism is the
linearity of curvature on a graph and the positivity-preserving
retraction; no higher-dimensional retraction statement follows.

